% Guide paragraphs for cote 134-6 sites 1--60, written after actual two-half collation.
\paragraph{92. Abelianization, its possible factorization, and universal chain models (pp. 320--331; 134-6/1--12).}
The entry dated 5.7 reorders the unfinished program at the end of part~IV. Questions about contractors, induced asphericity and contractibility, and some remaining material from part~II are to stay in this part; several existence and duality questions are deferred. The organizing issue is whether an arbitrary suitably qualified test category can do the work familiar from the simplicial, cubical, and hemispherical examples. The list distinguishes test-functor existence, closed model structures, independence of derived categories from the choice of indexing category, finite-type conditions, and boundary operations with Dold--Puppe relations (pp.~320--321). The last question motivates the section, but the author now doubts that it should be subordinated to cellular decompositions of spheres. The unexpectedly weak axioms already sufficient for a test category suggest seeking a more flexible construction of chain complexes. The additional ``strict'' qualification and the deletion of ``or contractor'' in the initial comparison are visible revisions to the scope of the claim.

The discussion fixes the ordinary localizer of weak homotopy equivalences and calls $\mathrm{D.}(\mathrm{Ab})$ by the name $\mathrm{Hotab}$ (pp.~321--322). Its starting point is the canonical abelianization functor $\mathrm{Hot}\to\mathrm{Hotab}$. Diagram (3) constructs it using a chosen test functor on $\Delta$, the free-abelian-group or ``Whitehead'' functor $\mathrm{Wh}_{\Delta}$, and the Dold--Puppe equivalence. The weak equivalences used on abelian objects are those detected on their underlying semisimplicial sets. Here ``Whitehead's theorem'' names the compatibility needed to descend free abelianization to the localized categories. Different choices of the test functor need not give isomorphic functors before localization; the claimed canonical independence afterward is referred to section~77, replacing a canceled reference to section~82 (pp.~322--323). The corresponding cubical and hemispherical comparisons are proposed with an express qualification about their Whitehead theorem. The following remark also separates the construction of abelianization itself from these stronger equivalence theorems: the ordinary unnormalized chain complex already sends weak equivalences to quasi-isomorphisms and yields the same localized functor.

For every small category $A$, formula (6) therefore gives a canonical functor $\mathrm{Hot}_A\to\mathrm{Hotab}$ through $\mathrm{Hot}$; no claim that $A$ is a test category is needed for that composition (p.~323). The new question is whether it factors through the localization $\mathrm{Hotab}_A$ of abelian presheaves, or can be computed by a functor from abelian presheaves to $\mathrm{Ch.}(\mathrm{Ab})$ (p.~324). These are additional requirements. An aspheric functor $j:\Delta\to A^{\wedge}$ gives a route back to the familiar simplicial case, but the text notes that its existence already forces substantial asphericity conditions on $A$. The comparison with a general $j:B\to A^{\wedge}$ leads to diagrams (11)--(14), including an independence question for the resulting factorization. In the contractor case, the argument is referred to the product remarks of section~82, p.~272 (p.~325). This reference has a different role from the corrected section~77 citation above. The live superscript on $j$ in (10') is not fully legible in the scan and remains marked as uncertain in the transcription.

Pages~326--328 attempt a construction without assuming such a $j$. Write $u=i^*i_A:A^{\wedge}\to\Delta^{\wedge}$. Although $u$ generally fails to preserve finite products, it preserves fiber products. Passing to the slice over $F=u(e_{A^{\wedge}})=i^*(A)$ gives an exact functor $u_o$ and hence a functor on abelian group objects. The localization morphism determined by $F$ supplies the left adjoint $\alpha^{ab}_!$ used in diagram (17). The lower square commutes by the adjunction argument, but the upper square is explicitly described as failing to look commutative. The proposed remedy is that its comparison map should become a weak equivalence after passage through the lower square. The author immediately revises the claim that this is the exact sufficient hypothesis: the composite functor (18) must also preserve weak equivalences so that it descends to localization. Neither condition is established for arbitrary $A$, and the text admits uncertainty even in the presence of the earlier aspheric functor. Diagram (19) then raises a further question: when $A$ models ordinary homotopy types, does the corresponding abelian comparison also become an equivalence? This is an open prospect in the section, not a conclusion of the diagrammatic construction.

The author interrupts this detour on p.~328 and returns to the desired computational description. Diagram (20) asks for a functor $L:A^{\wedge}_{ab}\to\mathrm{Ch.}(\mathrm{Ab})$ whose composite with free abelianization realizes the already defined homotopy abelianization. The distinction matters: the canonical functor to $\mathrm{Hotab}$ is available, while a particular chain-complex model and its factorization through abelian presheaves require further work. Using an aspheric simplicial comparison is called ``cheating'' because it imports the familiar model instead of producing the sought intrinsic construction. General lifting results into a category of models are mentioned as another possible route and then set aside (p.~329). The visible changes from ``abelian'' to ``additive'' for the derived category and from ``chain'' to ``abelian'' for its complexes sharpen the setting without supplying such a theorem.

The replacement idea is to vary the coefficient category. For any additive category $M$, a chain construction should give $L_M:\underline{\mathrm{Hom}}(A^o,M)\to\mathrm{Ch.}(M)$ compatibly with additive functors in $M$ (p.~329). The ``envelopping additive category'' $\mathrm{Add}(B)$ is introduced by a universal map from $B$: composition with it induces an equivalence between additive functors out of $\mathrm{Add}(B)$ and arbitrary functors out of $B$. Taking $B=A^o$, the text identifies such a compatible system with a chain complex $L.$ in $\mathrm{Add}(A^o)$. Formula (26), $L(X)=\widetilde X(L.)$, then evaluates this universal complex by the additive extension of an abelian presheaf $X$ (p.~330). The resulting existence problem concerns a complex whose induced square commutes; the associated uniqueness problem asks, for example, about uniqueness up to chain homotopy and about the category of pairs consisting of a complex and a compatibility isomorphism. These remain questions. The universal map is renumbered (23), while the sentence immediately after (24) still prints the reference (22); the source's uncorrected reference is retained.

The final revision restricts the apparent example available from a simplicial comparison (pp.~330--331). After first asserting that the construction obtained ``by cheating'' visibly comes from a universal complex, Grothendieck cancels an attempted justification and adds that $j$ must factor through $A$ itself. An aspheric functor $\Delta\to A$ then induces $\mathrm{Add}(j^o):\mathrm{Add}(\Delta^o)\to\mathrm{Add}(A^o)$, and the image of the familiar complex gives the desired candidate $L.$. Landing merely in $A^{\wedge}$ does not justify this particular construction. The section closes at that narrower sufficient situation. The next dated entry begins the promised explicit construction of the additive envelope, rather than resolving the wider existence and uniqueness questions.

\paragraph{93. Additive envelopes, cellular signs, and the proposed standard abelianizator (pp. 331--338; 134-6/12--19).}
The entry dated 6.7 makes the universal category from section~92 explicit. First replace each set of maps in $B$ by the free abelian group on those maps; then add formal finite families of objects, with morphisms given by matrices of such linear combinations (formulas (1)--(4), pp.~331--332). Matrix multiplication extends composition in $B$. The formal sums are deliberately kept distinct from any sums that already exist in $B$, since the desired universal property concerns arbitrary functors from $B$ into additive categories. The text states that the resulting $\mathrm{Add}(B)$ has the required two-dimensional universal property. A coefficient-ring variant replaces $\underline Z$ by a commutative ring $k$ with unit and targets additive categories equipped with $k\to\mathrm{End}(\mathrm{id}_M)$; the word ``commutative'' and the instruction to replace $\underline Z$ by $k$ are typed interlinear additions. This leads to the prospect of $k$-linearization and an associated basic localizer $\underline W_k$, without developing that theory here (p.~332).

For $B=A^o$, an \emph{abelianizator} is proposed as a pair $(L.,\lambda)$: a chain complex in $\mathrm{Add}(A^o)$ together with the compatibility isomorphism that makes section~92's square (20) compute homotopy abelianization (p.~332). The induced additive functor on abelian presheaves is available from this universal complex, but preservation of weak equivalences as quasi-isomorphisms is raised as an additional issue. The author considers making that preservation part of the definition if it is not automatic. Only after it holds does localization give the sought functor $\mathrm{Hotab}_A\to\mathrm{Hotab}$. The resulting square (19) of section~92 is then accounted for, while the earlier diagrams of type (15) are still described as less evident. Thus the terminology does not close the existence, uniqueness, or compatibility questions raised in section~92.

The cellular construction begins with the three standard categories and their finite products (pp.~333--334). For an object $a$, the ordered set $i(a)$ consists of its subobjects in $A$, explicitly distinguished from sieves on $a$. In the proposed setting it is finite, has a largest element, and gives the dimension of $a$ through the dimension of its geometric realization. The universal term is $L_n=\bigoplus_{\dim a=n}a$, requiring finite dimension fibers because the additive envelope has only formal finite sums. The differential on a summand $a$ is described by a linear combination of maps $b\to a$ from dimension $n-1$; in the standard cases its nonzero terms are precisely the boundary monomorphisms, with coefficients $\varepsilon_{\partial}\in\{\pm1\}$ (formulas (10)--(14)). The revision from ``any'' to ``all'' monomorphisms on p.~334 is therefore substantive. Choosing orientations of the cells $|i(a)|$ determines these coefficients by the stated Stokes rule. The text immediately identifies a difficulty: a subobject can have several representing monomorphisms if isomorphisms are allowed. Its first safe formulation assumes that every isomorphism is an identity.

Under those hypotheses, the anticommutativity of successive boundary orientations is invoked to obtain $d_{n-1}d_n=0$ (formula (18), p.~335). If nontrivial automorphisms are admitted, one must choose a representative monomorphism for each boundary subobject and impose the commutativity of the resulting diamonds (19). The page calls this a delicate condition, probably difficult to ensure. An additional factorization of every arrow as an epimorphism with section followed by a monomorphism would connect the construction with section~91 and make the subobject assignment a functor to ordered sets, but the author doubts that the existence of that functor is essential to the abelianizator question. There is also a narrower independence result: for different choices of orientations, the signs comparing them define a transitive system of componentwise isomorphisms of the complexes (22)--(25), p.~336. The circled marginal $\varepsilon_a^{\omega,\omega'}$ clarifies the sign notation. This identifies the orientation-dependent complexes canonically in the safe case. When representative boundary maps also have to be chosen, the text does not establish an analogous independence from those additional choices.

The final pages compare the scope of this construction with the test-category program (pp.~336--338). The familiar rigid examples and their finite products satisfy several desirable properties at once: the cellular condition, a canonical abelianizator, and strict test-category behavior. The author wonders whether these ``multistandard'' examples exhaust the possibilities, then contrasts unordered simplices, which retain automorphisms, with $\Delta^f$, which has no degeneracies and is expected here only to be a weak test category. Returning to the full subcategory $A_o$ of ordered sets considered in section~91, p.~318, the text argues that orientations need not be added to its objects or morphisms. Strict inclusions of ordered subsets provide the suggested coherent boundary maps. It accordingly claims that the unchanged category is a strict test category, while qualifying that it is presumably not a contractor. The earlier assumption of closure under finite products is then recognized as too restrictive for the standard examples themselves. A proposed substitute asks for compatible cellular subdivisions of products $|i(a)|\times|i(b)|$ by cells of the form $|i(c)|$. Whether this suffices even to make $A$ a weak test category, and whether the associated complex is an abelianizator, remain open questions in the entry.

The concluding remark revises its own terminology while it is being written (p.~338). ``Standard abelianizator'' is acceptable only if the constructed canonical complex actually has the required abelianization property in every case under consideration. ``Dold--Puppe chain complex'' is first offered as another name, with normalization viewed in the universal additive category. The text then corrects the attribution: the name is more appropriate for the \emph{normalized} complex obtained from $L.$, whereas the full complex itself is described as long familiar. Performing that normalization universally also requires enlarging the additive envelope to admit the direct summands associated with projectors. The two opposed marginal bracket arrows mark the end of this remark and the start of section~94; they do not provide a further proof or a resolved replacement for the open conditions above.

\paragraph{94. Recovering homotopy types from abelianization and varying the ground ring (pp. 338--342; 134-6/19--23).}
The emphasis changes from constructing abelianization through different modelizers to asking what extra structure would allow the homotopy type to be recovered from it. For a pointed connected type with fundamental group $\pi$, the first proposed ingredient is the chain complex of the universal covering, considered in $\mathrm{D.}(\underline Z(\pi))$ (formula (26), p.~339). This retains the group action needed for twisted coefficients. The second ingredient is a diagonal $L.\to L.\overset{L}{\otimes}L.$ in $\mathrm{Hotab}$, subject to suitable commutativity and associativity conditions (27). The text suggests combining these structures to include cup products with arbitrary twisted coefficients. Its larger question is whether they suffice to recover the entire homotopy type and therefore the other operations and invariants. That sufficiency is expressly a question which the author says he had not investigated.

The recollection of Sullivan's explanation then distinguishes such a diagonal, or the resulting rational cohomology ring under suitable finiteness assumptions, from a differential graded algebra carrying the fuller information (pp.~339--340). The author recalls a rational De Rham construction for spaces or semisimplicial sets and an equivalence after localization and restriction to the appropriate simply connected objects, but repeatedly qualifies his memory. The inserted words ``differential graded'' and ``in the derived category'' sharpen the statement: the relevant comparison is between algebra objects with their differential, up to the stated localized notion of isomorphism. The passage asks why such an algebra need not be recovered from its cohomology algebra with zero differential. It is not asserting that the cohomology ring alone classifies rational homotopy types, and the guide should not turn this remembered explanation into an unqualified theorem with unstated hypotheses.

Integral coefficients lead to the next proposal (pp.~340--341). The reported obstruction to using an anticommutative differential graded algebra prompts comparison with the failure of the algebraic Poincar\'e lemma and its divided-power remedy. Grothendieck recalls developing a De Rham formalism with divided powers for arbitrary semisimplicial sets, and considers its dual coalgebra form technically more appropriate. He wonders whether a suitable derived category of this structure could recover homotopy types. The later restatement specifies a bigraduation, and for nontrivial fundamental group it again invokes the universal covering and its group action. These are proposed structures and recollections of earlier work; no reconstruction theorem, equivalence relation on all such objects, or proof of sufficiency is supplied here. The discussion of derived categories and binomial coefficients explains what attracted his attention, rather than completing that missing argument.

The feature that drives the remainder of the section is the possibility of defining the algebraic structure over an arbitrary commutative ring, or even a scheme (pp.~341--342). This suggests a notion of homotopy types over a ground ring and a change-of-ring operation. Abelianization provides the familiar model through the dependence of $\mathrm{D.}(k)$ on $k$, but the proposed nonabelian version is less settled. The author recalls expressing some simple spaces, including standard $K(\pi,n)$ examples and fibrations between them, through integral points of finite-type affine semisimplicial schemes. A prospective functor from $\mathrm{Hot}$ to a suitable derived category of such schemes would make naive scheme-theoretic base change $\underline Z\to k$ reflect the expected change of coefficients. The text entertains an equivalence, then stresses that the attempt was brief, was not resumed, and may have been unrealistic. These qualifications belong to the mathematical status of the proposal.

The final paragraph tries to extend that picture beyond simply connected examples by using nilpotent fundamental groups (p.~342). It first hopes to recover a suitably restricted finitely presented torsion-free group as the integral points of a finite-type affine nilpotent group scheme, then retreats to a rational group scheme and recovery only up to commensurability. The handwritten ``Possibly'' introduces an anticipated equivalence after localizing the discrete groups at monomorphisms of finite index. On the algebraic side the page says connected nilpotent affine group schemes over $\underline Q$, followed by the further description that their underlying schemes are affine spaces. That latter restriction must remain visible when interpreting the proposed class; the paragraph does not establish the stated equivalence of descriptions or a categorical equivalence. The next section leaves these afterthoughts and returns to the contractor program of section~92.

\paragraph{95. Contractors and the comparison of contractibility with asphericity (pp. 343--351; 134-6/24--33).}
The entry dated 7.7 resumes the three topics announced in section~92, beginning with contractors. Its first definition is canceled because it tacitly requires finite products in $A$, an assumption which excludes the standard test categories that the author wants to include (p.~343). The replacement embeds $A$ in $A^{\wedge}$ and asks that the intervals coming from its objects generate a contractibility structure in which those objects are contractible. This is the condition called Bas~4 in section~51, p.~118. A category satisfying it is first called a \emph{precontractor}; further compatibility with asphericity will define a contractor. The inserted qualification that the ambient category $M$ is stable under finite products belongs to the generalized formulation of this condition. The discarded initial definition and the separate abandoned version of p.~345 remain visible stages of the argument, rather than alternative final definitions.

Proposition~1 explains how the structure passes from $A$ to another category $M$ stable under finite products (pp.~344--345). The three sufficient situations are: $A$ itself has finite products and the functor preserves them; an extension from $A^{\wedge}$ exists and preserves the final object and binary products of representables; or the original functor is fully faithful and strictly generating. The discussion relates these conditions to the criteria on pp.~281 and 283 and to full faithfulness of the associated $f^*:M\to A^{\wedge}$. In the last case, identifying $M$ with a full subcategory containing $A$ gives a direct argument without separately requiring all small inductive limits or the existence of $f_!$. The proposition's universal validity in the indicated cases also characterizes the precontractor condition. These distinctions explain why the author replaces an intrinsic product assumption on $A$ by a condition on its presheaf category.

Proposition~2 supplies the missing asphericity requirement (p.~345). For a precontractor, let $h$ be generated by the intervals coming from $A$, and compare its contractible objects with the canonical asphericity structure for ordinary weak equivalences. The listed equivalent conditions include total asphericity of $A^{\wedge}$, generation of that asphericity structure by the contractibility structure, the corresponding inclusions of contractible objects, compatibility of $h$-homotopisms with weak equivalences, and equality with the homotopy structure associated with those weak equivalences. The pen changes from ``totally'' to ``locally'' in the object-level conditions are consequential: the initial notation $A^{\wedge}_{\mathrm{totas}}$ is not consistently replaced everywhere, but the revised conditions and Definition~1 on p.~346 use local asphericity for the representables. The resulting \emph{contractor} is a precontractor satisfying these equivalent conditions. The page then records trivial examples, generating subcategories of contractibility structures, and the three standard categories with their finite products.

The revision of Proposition~3 separates the existence of a separating interval from the additional asphericity needed for a strict test category (pp.~346--347). A nontrivial precontractor contains such an interval; when it is also totally aspheric, hence a contractor in the terminology here, the strict test-category conclusion follows. The proof uses sections of objects to show that an interval is either separating or has equal endpoints. If every interval from $A$ had equal endpoints, the resulting homotopy relation would collapse to equality and every contractible object would be final, contradicting nontriviality. The subsequent corollary characterizes a trivial contractibility structure by several equivalent collapses; a fifth handwritten condition says that any two sections of a contractible object are equal. The supplied scan cuts off faint traces below that addition, and the input transcription supplies no hidden continuation there. The guide therefore does not reconstruct one.

For a nontrivial contractibility structure $(M,M_c)$, the next argument transports a contractible interval with distinct endpoints through $i^*:M\to A^{\wedge}$ and obtains a separating interval (p.~348). With the stated total $\underline W$-asphericity assumptions this yields a $\underline W$-test category, and the stronger assumptions give a strict one. The author explicitly replaces corollary~3 on p.~253: the earlier multiplicative-interval and membership assumptions are declared unnecessary, and the separation requirement reduces to distinct endpoints. The following remark corrects a different possible confusion. The structure $h$ generated by intervals from $A$ must be distinguished from $h'=h_{\underline W_A}$ used in section~79; the latter can depend on the localizer. The diagram's lower adjoint index remains visually uncertain between a heavily retraced shriek and a star, so the transcription does not resolve it from the adjunction argument.

Under strict total $\underline W$-asphericity, together with Loc~4, the text then compares both structures with the canonical homotopy structure defined by 0-connected intervals (pp.~349--350). The coarsest localizer under discussion is described by bijectivity on $\pi_o$. The proposed new auxiliary proposition on p.~349 is crossed out and replaced by a reference to Lemma~2 on p.~268; its blank topos names are left blank in the source. The surviving Proposition~4 states the equality $h=h_{\underline W_A}=h_{A^{\wedge}}$, then equates homotopy compatibility of $i^*$, $\underline W$-asphericity of $i$, and the specified 0-connectedness criterion on a generating subcategory of $M_c$. Nontriviality of the contractibility structure gives the strict $\underline W$-test conclusion. The proof on p.~351 refers back to section~79, p.~252, and Proposition~3's corollary. If $M_c$ is moreover $\underline W$-modelizing, the final corollary makes $i$ a $\underline W$-test functor and gives the equivalence between the displayed localizations. The next dated entry returns to the definitions themselves, particularly the treatment of the empty category.

\paragraph{96. Verticality, horizontal topoi, and the limits of the terminology (pp. 351--354; 134-6/33--36).}
The entry dated 8.7 begins by explaining the repeated abandoned pages in section~95: the author had expected a smooth exposition but found the technical adjustments less settled than anticipated. The first terminological issue separates two properties of a topos: every nonempty sheaf has a global section, and the final object has only the empty and full subobjects (pp.~351--352). The typed replacement of ``as being equivalent'' by ``to imply'' is mathematically significant. The classifying topos $B_G$ is then used to distinguish the properties for a nontrivial discrete group. The first property is discussed as a particularly strong global asphericity condition, initially only as far as $H^1$ is concerned; the inserted ``global'' specifies the intended cohomological scope. Several possible names for the second property are considered before the author settles on ``vertical,'' with the image of a small base and a tower of internal structure.

The explicit definition excludes the empty topos and requires its open subtopoi to be only the empty and full ones (p.~352). A horizontal topos is defined by generation of its category of sheaves by subobjects of the final object; the topos of a topological space is the principal example (p.~353). The revised statement says that being both horizontal and vertical gives a final topos, represented by the one-point space and the category of sets. The preceding discussion's accidental inclusion of the empty case is thereby removed. The author also notices that the initial precontractor definition formally allowed the empty category, which he does not want to include.

There is a scope problem in the accompanying equivalences which the guide must leave explicit. Pages~351 and 353 identify the vertical condition for a topos of the form $A^{\wedge}$ with the requirement that every object of $A$ have a section, and p.~354 repeats the claimed coincidence. If $A$ is an unrestricted small category, the source's own $B_G$ example supplies a counterexample: take $A$ to be a nontrivial one-object group category. Its presheaf topos has only the empty and full subterminal objects, but the regular $G$-set has no fixed point and hence no global section. Thus the converse implication requires an additional restriction on $A$, or a correction to the assertion. The pages do not state that restriction at these occurrences. The transcription preserves the claims; the guide does not turn them into an unrestricted equivalence.

Verticality is nevertheless presented as a separate direction from total asphericity or total 0-connectedness, with Lemma~2 on p.~268 and section~60, p.~149 as motivations (p.~353). Grothendieck therefore proposes replacing the earlier phrase ``strictly totally aspheric'' by ``totally aspheric and vertical,'' rather than treating the extra condition as a degree of asphericity. A second terminological conflict remains unresolved: the notions of locally and totally aspheric objects introduced in section~79 do not automatically agree with the terminology obtained by passing to the associated topoi, as the principle of section~66 would suggest. The author elects to keep this useful but unorthodox language until it causes a concrete difficulty. The section closes by returning to the three structures $h,h',h''$ from section~95: they coincide for a contractor, while only $h'$ depends on the chosen basic localizer. Clarifying their general relationship is announced as the next task, then interrupted before the entry dated 12.7.


\paragraph{97. Projective topoi, the three homotopy structures, and two notions of morphism (pp. 354--357; 134-6/36--39).}
The entry dated 12.7 separates the two conditions which the previous discussion had partially conflated. A topos is called \emph{projective} when every epimorphism $F\to e$ to its final object has a section, equivalently when that final object is projective in the category of sheaves (pp.~354--355). Among nonempty topoi, the property that every nonempty sheaf has a section is then the conjunction of projectivity and verticality. The cohomological discussion uses dimension shifting to pass from vanishing of $H^1$ with arbitrary abelian coefficients to vanishing of all higher cohomology; the pen addition ``except for $H^0$'' and the typed word ``global'' delimit the assertion. A further handwritten qualification requires 0-connectedness in the claimed asphericity consequence. The page again asserts that the two original conditions coincide for topoi $A^{\wedge}$, so its inference for arbitrary vertical categories inherits the scope problem identified in section~96. The distinction between projectivity and verticality itself remains explicit and useful.

The promised comparison diagram on p.~355 identifies $h=h_{A^{\wedge},A}$, generated by intervals coming from $A$; $h'=h_{\underline W_A}$, depending on the localizer; and the canonical structure $h''=h_{A^{\wedge}}=h_{\underline W_{oA}}$. It displays three \emph{conditional} inequalities $h\leq h'\leq h''\leq h$. Their separate hypotheses are total $\underline W$-asphericity of $A$, the inclusion $\underline W\subset\underline W_o$, and verticality of $A$. The inserted words ``conditional'' and ``on A or W'' guard against reading this as an unconditional cycle. The first hypothesis also bears a circled marginal ``\'ecarter'' connected to its bracket. That mark does not securely establish a deletion, so the diplomatic text preserves the hypothesis, connector, and instruction together. The revised sentence below the diagram states the coincidence of all three structures under the combined assumptions; taking $\underline W=\underline W_o$ recovers the comparison with Lemma~2 on p.~268. The following sentence refers to (1), although the displayed comparison is labelled with a star. This source cross-reference is left unchanged.

For a contractor, the three structures are accordingly said to coincide for every localizer satisfying the present Loc~4 (pp.~355--356). If the contractor is nontrivial, the strict test-category and modelizing results give a \emph{canonical modelizer}: the weak-equivalence structure is determined by the canonical homotopy structure. Alongside $(\mathrm{Cat})$ and $(\mathrm{Spaces})$, these presheaf categories are the principal examples available in the text. Stacks are proposed as another substantial family, with no construction or proof supplied in this entry. The point of the comparison is therefore to identify when the modelizing structure depends only on the canonical homotopy relation, rather than on a separate choice of generating intervals.

The first proposed definition of a morphism of contractors $f:A\to B$ requires $f^*:B^{\wedge}\to A^{\wedge}$ to preserve the homotopy structures (p.~356). Conditions (2)--(5) express this through preservation of contractible objects, through the representables in $B$, through their 0-connectedness, or through their $\underline W$-asphericity. The A/B exchanges in these displays are visible pen corrections and are essential to the direction of the conditions. The criterion reduces to $\underline W$-asphericity of $f$, and the stronger-looking equality (5') identifies the entire asphericity structure of $B^{\wedge}$ as the inverse image under $f^*$ of that on $A^{\wedge}$. The author recognizes this as an instance of Proposition~4 of section~95, p.~350, after initially presenting it as a fresh surprise.

The final page tests a broader situation in which $f$ lands in $B^{\wedge}_c$, without necessarily factoring through $B$ (p.~357). The associated adjoint pair $f_!\dashv f^*$ suggests the two requirements (8), one on $f_!(A)$ and one on $f^*(B)$. Closure under composition leads to the stronger condition (9), $f_!(A^{\wedge}_c)\subset B^{\wedge}_c$, which follows when $f_!$ preserves finite products by the corrected reference to section~85. It need not follow merely from (8), even when the original functor lands in $B$. The author therefore distinguishes two plausible notions, neither claimed to imply the other: he reserves ``morphism of contractors'' for the first and calls the second a bimorphism of the associated contractibility structures, referring to section~86. The entry closes without deciding which will be more useful. The dated section~98 then postpones the topic of induced structures and turns to simplices without degeneracies.

\paragraph{98. Simplices without degeneracies and the remaining asphericity lemma (pp. 357--360; 134-6/39--42).}
The entry dated 16.7 postpones induced asphericity and contractibility structures on slice categories and returns to the assertion announced in section~43: the category $\Delta^f$ of ordered standard simplices with injective maps should be a weak test category. The author recalls having found a simple proof in March without keeping notes, and reports difficulty recovering it (pp.~357--358). Diagram (1) compares ordered and unordered simplices, with and without degeneracies. The ordinary categories $\Delta$ and $\widetilde\Delta$ are contractors; their injective-map variants are said not to be test categories. The claim pursued here is the weaker one for $\Delta^f$, with the analogous result for $\widetilde\Delta^f$ suggested as a possibility. The diagram's lower arrow is literally labelled $\beta f$, and the blank before an isolated subscript $n$ in the following sentence remains unfilled in the scan.

The expected common property of all four comparison functors is asphericity, but only $\alpha$ and $\beta$ are treated as checked for the immediate argument (p.~358). The unordered-simplex functor $\beta$ is moreover said to be a morphism of contractors: choosing a point in a finite nonempty set $E$ gives an elementary homotopy from the identity of $\beta^*(E)$ to the associated constant map. The author does not claim that $\beta$ is a bimorphism of the canonical contractibility structures. That stronger question would follow from preservation of finite products by $\beta_!$, which remains unproved here. This is the first concrete use of the distinction made at the end of section~97.

The proof strategy identifies the category of injective simplices over $\Delta_n$ with the ordered set of its nonempty subsets (p.~359). Consequently $i_{\Delta^f}$ is canonically isomorphic to the composite of $\beta\alpha$ with the standard unordered-simplex test functor $\widetilde i$ from section~34. The known asphericity of $\widetilde i$, combined with asphericity of the comparison, would yield the required result. The page explicitly fixes the strongest localizer under discussion, $\underline W_{\infty}$ of ordinary weak equivalences. To prove asphericity of $\alpha$, it adjoins an initial empty simplex to form $\Delta'^f$ and identifies each relevant slice with $(\Delta'^f)^{n+1}\setminus U^{n+1}$, where $U$ is the open subcategory consisting of that initial object. The inserted U is typed; the marginal handwritten formula and the canceled exponent mark show the local revision of this description.

Two lemmas remain in that reduction. Lemma~1 asserts the asphericity of $\Delta^f$ itself. Lemma~2 says that, for a finite \emph{nonempty} family of pairs consisting of a category and an open subcategory, asphericity of each ambient category and each closed complement implies asphericity of the complement of the product of the opens (pp.~359--360). The nonemptiness qualification is a handwritten insertion. The proof reduces to two factors and applies a Mayer--Vietoris argument to the two closed pieces and their intersection. The resulting set difference $X\setminus U$ is visible in the source and must not be changed into a product.

The text then acknowledges that Lemma~1 is where the proof has stalled (p.~360). A heuristic calculation would use the familiar boundary operations for semisimplicial complexes, assuming that the same cohomological formulas apply to face complexes. Applied to the final face complex, and including noncommutative $H^1$ with constant coefficients, this would yield the required asphericity and even establish that of $\alpha$ directly. The entry does not supply that calculation or justify the extension of the formulas at this point. Instead, trying to understand those boundary operations brings the author back to sections~92--93 and to a stronger notion of abelianizator with a prospective existence and uniqueness theorem. Section~99 begins that reformulation before the unfinished asphericity argument is completed.

\paragraph{99. Integrators, arbitrary coefficients, and the finite-sum restriction (pp. 360--366; 134-6/42--48).}
The uneasiness left by sections~92--93 is now traced to their concentration on constant coefficients. Grothendieck's habitual cohomological viewpoint would instead begin with arbitrary sheaves and regard constant or locally constant coefficients as particular cases (pp.~360--361). He therefore puts $B=A^o$ and asks for another interpretation of the same universal chain complex: it should compute cohomology of $B$, meaning its presheaf topos, with coefficients in an arbitrary abelian presheaf on $B$. The variance is essential. The objects whose homotopy types motivated the earlier construction lie in $A^{\wedge}$; the coefficient objects in this reformulation lie in $B^{\wedge}_{ab}$. The inserted ``abelian'' is typed, while the circled marginal $A^{\wedge}_{ab}$ and the inserted phrase locating the earlier abelianizations keep the two sides of this comparison explicit.

The additive envelope is first given a concrete realization (pp.~362--363). The canonical inclusion $B\to B^{\wedge}$ followed by free abelianization factors through an additive functor $\gamma_B:\mathrm{Add}(B)\to B^{\wedge}_{ab}$. The universal property identifies abelian presheaves $F$ with additive functors $\widetilde F:\mathrm{Add}(B)^o\to(\mathrm{Ab})$, and formula (5) identifies evaluation on $L$ with $\mathrm{Hom}(\gamma_B(L),F)$. This makes $\gamma_B$ fully faithful and realizes its image as the finite direct sums of free abelian representables. The source's formula (5) prints only $B^{\wedge}$ in the Hom index; p.~364 explicitly specifies that the Hom is taken in the category of abelian presheaves. Likewise, the generalization (3') omits an additive qualifier on its right-hand Hom even though the prose requires additive functors. The guide follows those explicit descriptions while the transcription preserves the displayed notation. The revision from ``can be generalized into'' to ``is a particular case of,'' and the typed replacement of ``abelian'' by ``additive,'' clarify the intended universal-property argument.

The image of $\mathrm{Add}(B)$ consists of projectives of finite presentation, but it need not contain every such projective (p.~363). The exact statement is that every projective of finite presentation is a direct factor of one of these finite sums. The typed replacement of ``type'' by ``presentation'' and the insertion ``full sub'' sharpen this description. Passing to the Karoubi envelope, which formally supplies images of projectors, gives the full subcategory of projectives of finite presentation. Its proposed universal property targets additive categories in which every projector has an image. This distinction between finite sums and their retracts will matter when the author asks whether the components of a resolution can lie in the original additive envelope.

For a chain complex $L.$ in that envelope, formula (6) identifies $\widetilde F(L.)$ with the componentwise Hom cochain complex (p.~364). The desired computational property is a natural derived isomorphism $\mathrm{Hom}^{\bullet}(L.,F)\simeq R\Gamma_B(F)$ for every abelian presheaf $F$. The text interprets $R\Gamma_B(F)$ as $R\mathrm{Hom}(\underline Z_B,F)$, where $\underline Z_B$ is the constant presheaf, and uses projectivity of the components to compare ordinary Hom with derived Hom. In the projective-resolution setting this leads to the augmentation $L.\to\underline Z_B$ and the requirement that it be a quasi-isomorphism. The page's argument is the familiar one for resolutions in nonnegative chain degrees; it supplies no separate justification for an arbitrary unbounded complex of projectives. The lower index of the projective limit in (7) is visibly corrected to $B^{\circ}$, with a circled marginal repetition, while the older subscript remains partly unreadable.

An \emph{integrator} on $B$ is consequently defined to be a projective resolution of $\underline Z_B$ (p.~365). It is \emph{special} when its components belong to $\mathrm{Add}(B)$. Since $\underline Z_B$ need not itself belong to that category, the augmentation is extra structure external to the additive envelope; formula (11) encodes it by the element $\lambda$ obtained from evaluation on $L_o$. The usual comparison theorem gives chain-homotopy uniqueness of integrators, and fullness gives the corresponding statement within $\mathrm{Add}(B)$ for special ones. General integrators exist because the abelian presheaf category has enough projectives. Existence of a special integrator is a further finiteness question: successive kernels need not be quotients of finite sums of representables. The page later says ``strict integrator'' once where the discussion has concerned the defined term ``special''; that wording is retained in the source.

The one-object group category supplies the proposed test case. A special integrator would be a resolution of the trivial $\underline Z[G]$-module $\underline Z$ by free modules of finite rank in each degree. The author doubts that this can happen unless $G$ is finite, and a marginal X stands beside that sentence (p.~365). The mark does not establish a deletion or a revised theorem. The proposed necessity of finite $G$ is nevertheless false: for the infinite cyclic group, write $R=\underline Z[t,t^{-1}]$ and use the exact augmented sequence $0\to R\xrightarrow{t-1}R\to\underline Z\to0$. Multiplication by $t-1$ is injective, and its cokernel is the augmentation module, so both nonzero resolution terms are free of rank one. This example leaves intact the author's broader observation that special integrators express a homological finiteness condition, while preventing the guide from adopting his tentative finite-group criterion.

The immediate construction relaxes the finite-sum requirement (pp.~365--366). The category $\mathrm{Addinf}(B)$ is formed by allowing arbitrary small direct sums and has the corresponding universal property for functors into infinitely additive categories. It embeds in $B^{\wedge}_{ab}$, its objects remain projective, and every abelian presheaf is a quotient of one of them. Its retracts therefore give all projectives; the Karoubi-envelope formulation is repeated with the appropriate infinite-additivity condition. An integrator with components in $\mathrm{Addinf}(B)$ is called \emph{quasi-special}. Such resolutions always exist, and the text asserts their uniqueness up to chain homotopy inside this larger category. The typed insertion ``up to equivalence'' qualifies the categorical identification; its long placement bracket passes across the following words without deleting them.

The last sentence extends the description of augmentation data in (11) to the quasi-special case, allowing the indexing set $I_o$ to be infinite (p.~366). A notation issue must then be kept in view. For $L_o=\bigoplus_{i\in I_o}\mathrm{Wh}_B(b_i)$, a map $L_o\to\underline Z_B$ is specified by an arbitrary family of integers, hence by an element of $\prod_{i\in I_o}\underline Z$. Formula (11) prints $\underline Z^{(I_o)}$, which agrees with the product for finite $I_o$. If its parenthesized exponent denotes the direct sum, the final extension needs that notation changed or interpreted as a product; an unrestricted augmentation need not have finite support. The transcription leaves the formula and the extension as written. The next entry adopts a canonical notation for the homotopy class of the augmented resolution and develops its homological meaning.


\paragraph{100. Intrinsic homology, relative integration, and the comparison with cohomology (pp. 367--374; 134-6/49--56).}
The entry actually dated 17.6 adopts the notation $L^B_{\bullet}\to\underline Z_B$ for an integrator, unique up to homotopy as an augmented complex; the date is retained despite the surrounding July entries (p.~367). Formula (2) identifies $R\Gamma_B(F)$ with $\mathrm{Hom}^{\bullet}(L^B_{\bullet},F)$, and (3) gives the resulting cohomology groups. A quasi-special integrator has components which are arbitrary small direct sums of free abelian representables, while a special one has finite sums in every degree. The disjointness convention for the index sets lets the author suppress the degree on the objects $b_i$. A substantive qualification is needed in (6), p.~368: the scan has a clear handwritten direct-sum sign in the expression for $\mathrm{Hom}(L_n^B,F)$. In general this Hom is the product $\prod_{i\in I_n}F(b_i)$, because a map from a direct sum is an arbitrary family of maps from its summands. The displayed sum is therefore justified for the finite indexing sets of special integrators, but does not give the unrestricted quasi-special formula. This is the same finite-versus-infinite distinction left implicit in the augmentation notation at the end of section~99. The source sign remains unchanged. Formula (7) extends the construction to bounded-below complexes of coefficients by the simple complex associated with the componentwise double Hom complex. The specialization to a constant complex $K_B^{\bullet}=p_B^*K^{\bullet}$ uses the adjunction $p_{B!}^{ab}\dashv p_B^*$ to obtain (10), where Hom is now taken in abelian groups (pp.~368--369). The source's reference to (2') is retained even though the immediately relevant generalization is labelled (7). The resulting complex $p_{B!}^{ab}(L^B_{\bullet})$, consisting of projective and hence free abelian groups, is proposed as an intrinsic representative of the global homology of $B$ and of its homotopy type. Its covariance and the functoriality of (10) are invoked to descend the construction from categories to homotopy types, producing the abelianization functor (12). This should agree with the functor of section~92, p.~321, obtained by a particular test category and the diagram on p.~322. The proposed comparison would establish the same constant-coefficient formula and recover an object of $D^-(\mathrm{Ab})$ from its Hom functor on $D^+(\mathrm{Ab})$. The author explicitly has not checked the latter assertion, particularly its suggested generalization to arbitrary abelian categories; the guide therefore records the route to the comparison rather than presenting the sketch as a completed uniqueness proof. For a functor $f:B'\to B$, formula (14) introduces $L^{B'/B}_{\bullet}=f_!^{ab}(L^{B'}_{\bullet})$ as a relative homology object with projective components (p.~370). It remains quasi-special or special when the chosen source resolution is so. The text asserts that coasphericity of $f$, meaning asphericity of the opposite functor, makes this object an integrator on $B$; its proposed converse is qualified by the choice of the homological localizer $\underline W_{\infty}^{Z}$ and deferred to later notes. The first symbol and the raised correction in (14) retain local overprint doubts, but the prose and surrounding formulas make the intended relative construction explicit. The long first remark contrasts the author's experience with cohomology on general topoi and the more direct geometric history of homology (pp.~370--373). Universal-coefficient formulas, twisted coefficients on a topological variety, and the older sheaf-theoretic comparison results motivate the question, without supplying a single equivalence that preserves every coefficient system and extra structure. The text attributes the asymmetry in general topoi to the lack of enough projectives, despite the availability of enough injectives for abelian sheaves. Presheaf topoi $B^{\wedge}$ provide the favorable case with enough projectives. For morphisms induced by a functor between the indexing categories, the adjoint triples $f_!\dashv f^*\dashv f_*$ and $f_!^{ab}\dashv f^*_{ab}\dashv f_*^{ab}$ explain the additional left-derived operation (p.~372). The set-valued $f_!$ need not take group objects to group objects, so the abelian left adjoint is constructed by its own adjunction. The author proposes calling $Lf_!^{ab}$ \emph{integration} and $Rf_*$ \emph{cointegration}, revising his former use of the first name for the latter functor (p.~373). Projection from a disjoint family of copies of $B$ gives the motivating direct-sum versus direct-product picture. The inserted ``presumably'' qualifies the assertion that integration is as meaningful as cointegration, while the long pen addition says that the author's continuing personal emphasis on cohomology no longer amounts to confidence in its technical superiority. Possible dualities between homology for $f$ and cohomology for $f^o$ remain suggestions. The second remark allows a commutative coefficient ring $k$, replacing abelian groups and additive envelopes by $k$-modules and their $k$-linear counterparts; the exposition stays over $\underline Z$ only for simplicity (pp.~373--374). The next dated entry returns to the postponed computational description of $f_!^{ab}$ and the relation between integrators and abelianizators.

\paragraph{101. Computing integration and the limits of internal adjunction (pp. 374--377; 134-6/56--59).}
The computational description postponed in section~100 begins with the fact that $f_!^{ab}$ preserves small inductive limits and is right exact (p.~374). Every abelian presheaf on $B'$ has a presentation by two objects of $\mathrm{Addinf}(B')$, so the functor is determined by its restriction to this category of special projectives. Diagram (1) identifies that restriction with the tautological infinitely additive extension $\mathrm{Addinf}(f)$. Formula (2), p.~375, makes the description explicit: apply $f$ to the indexing objects of a direct sum of free abelian representables and take the corresponding direct sum on $B$. The free-abelianization identity (4) and diagram (5) explain why this gives $f_!^{ab}$; the page refers to SGA~4, I,5.8.3 for the analogous construction with other algebraic structures. Restriction to finite sums gives $\mathrm{Add}(f)$ as well. Diagram (6) collects the original categories, their finite and infinite additive envelopes, their presheaf categories and free abelianization into one commuting configuration. The circled instruction to center and spread it out and the unattached outer pencil curve are graphic annotations, not additional functors. The scan also shows that the diagonal from $B$ ends at the lower $B^{\wedge}$, a destination that must be preserved when the diagram is reconstructed. The passage to derived functors uses bounded-above complexes of these projectives (p.~376). Two visible corrections are decisive: ``below'' becomes ``above'' in the degree bound, and ``quasi-isomorphisms'' becomes ``homotopy equivalences'' in the localization description (8). On bounded-above complexes of projectives the two classes coincide, so the corrected formulation describes the usual homotopy model of $D^-(B^{\wedge}_{ab})$. Applying $\mathrm{Addinf}(f)$ componentwise then computes $Lf_!^{ab}$. The displayed right side of (9) literally has $\mathrm{Addinf}(L'_{\bullet})$, without the $f$ argument; the immediately following prose supplies the intended operation $\mathrm{Addinf}(f)$. The transcription leaves that formula as written, and preserves the separate prime and lower dot on the argument, repeated in the circled margin. Formula (10) applies the construction to the relative integrator $L^{B'/B}_{\bullet}$, while projection to the final category gives an explicit global homology complex whose $n$th term is $\underline Z^{(I_n)}$ (11). Here the parenthesized exponent really describes the direct sum: the projection sends each free abelian representable to $\underline Z$. The final comparison returns to cointegration with coefficients pulled back from $B$ and asserts the internal-Hom isomorphisms (12), (12') and (13), p.~377. These are presented as a localization of ordinary derived adjunction, with an arbitrary bounded-below coefficient complex $K^{\bullet}$ and, in (13), an arbitrary bounded-above $F'_{\bullet}$. An additional hypothesis is required for this internalization. A concrete counterexample already occurs in degree zero: let $B=[0\to1]$, let $B'$ be the final category, and let $f$ select $1$. An abelian presheaf on $B$ is a diagram $K_1\to K_0$. Then $f^*$ is evaluation at $1$, $f_!^{ab}(V)=(V\xrightarrow{1}V)$, and $f_*(V)=(V\to0)$; all three functors are exact. Taking the integrator $\underline Z$ on $B'$ gives $L^{B'/B}=f_!^{ab}\underline Z=(\underline Z\xrightarrow{1}\underline Z)$, the tensor unit and the free abelian representable associated with $1$. Internal Hom from that unit is the identity functor. For $K=(0\to\underline Z)$, therefore, the left side of (12) or (12') is $Rf_*f^*K=0$, while the claimed right side is $R\underline{\mathrm{Hom}}(L^{B'/B},K)=K\ne0$. The ordinary external adjunction remains valid: $\mathrm{Hom}(L^{B'/B},K)=K_1=0$. Equivalently, the projection-formula map needed here cannot be an isomorphism, since $f_!^{ab}(\underline Z\otimes f^*K)=0$ whereas $f_!^{ab}\underline Z\otimes K=K$. Thus existence of the three adjoints and the projective-resolution construction does not by itself justify the claimed internal formula; a suitable local or projection-formula compatibility must be added. All three source formulas remain unchanged, including the visibly missing outer closing parenthesis on the left of (13). The last remark calls this the main relation between homology and cohomology, or abelian integration and cointegration, and proposes the same formalism for ringed topoi possessing $f_!$ on sheaves of sets. Its cancellation of ``commutative'' explicitly permits more general sheaves of rings, while ``striking'' replaces ``interesting'' in the closing assessment of possible applications. That proposed generalization inherits the preceding qualification about internal adjunction; the page does not state the missing local condition. Section~102 next turns from these operations to the pairing between the mutually dual indexing categories and their abelian presheaves.
